\section{\label{sec:su2_viz} Connection to SU(2) and the Bloch vector}

The plane-dihedral construction of Section~\ref{sec:hyperchords} realizes, in geometric terms, the standard parametrization of a two-component spinor and the associated Bloch vector map.
Throughout this section we keep the scale convention
\[
u := \ssp^\dagger\ssp,
\]
and we recall the geometric embedding convention introduced in Section~\ref{sec:hyperchords_setup}:
the rendered scene is scaled by the rescaled spinor
\[
\widetilde{\ssp}:=\|\ssp\|\,\ssp=\sqrt{u}\,\ssp,
\qquad\text{so that}\qquad
\|\widetilde{\ssp}\|=u.
\]
This is why the visualization sphere is $S_u=\{\xx\in\Rt:\|\xx\|=u\}$.
All direction-only quantities (e.g.\ the Bloch vector) are invariant under this scaling.


\subsection{Standard half-angle parametrization}
\label{sec:standard_param}

Let $\ssp\in\mathbb{C}^2\setminus\{0\}$ and write $\rho:=\|\ssp\|=\sqrt{u}$.
Using the symmetric (Hopf-fiber) phase split adopted throughout this work, $\ssp$ can be written as
\begin{equation}
\label{eq:spinor_symmetric_halfangle}
\ssp
=
\rho\,e^{\iu\gamma}
\begin{bmatrix}
e^{-\iu\phi/2}\cos(\theta/2)\\[2pt]
e^{+\iu\phi/2}\sin(\theta/2)
\end{bmatrix},
\qquad
\theta\in[0,\pi],\;\phi\in[0,2\pi),\;\gamma\in[0,2\pi).
\end{equation}
Equivalently, in terms of component arguments
\[
\spu=|\spu|e^{\iu\psi_\uparrow},\qquad
\spd=|\spd|e^{\iu\psi_\downarrow},
\]
one has
\begin{equation}
\label{eq:phase_split_defs_secIV}
\phi=\psi_\downarrow-\psi_\uparrow\ (\mathrm{mod}\ 2\pi),
\qquad
\gamma=\tfrac12(\psi_\uparrow+\psi_\downarrow)\ (\mathrm{mod}\ 2\pi),
\qquad
\psi_\uparrow=\gamma-\phi/2,\ \psi_\downarrow=\gamma+\phi/2.
\end{equation}
The moduli are
\[
|\spu|=\rho\cos(\theta/2),
\qquad
|\spd|=\rho\sin(\theta/2),
\qquad
\frac{|\spu|^2}{u}=\cos^2(\theta/2),
\qquad
\frac{|\spd|^2}{u}=\sin^2(\theta/2).
\]

\paragraph{Gauge-fixed form (common in the literature).}
A frequently used gauge chooses a representative in which the upper component is real and nonnegative, yielding
\[
\ssp
=
\rho\,e^{\iu\alpha}
\begin{bmatrix}
\cos(\theta/2)\\
e^{\iu\phi}\sin(\theta/2)
\end{bmatrix}
\quad\text{for some }\alpha\in[0,2\pi),
\]
which is equivalent to \eqref{eq:spinor_symmetric_halfangle} after absorbing a phase into $\alpha$.
We use the symmetric split \eqref{eq:spinor_symmetric_halfangle} because it separates the relative phase $\phi$ from the common (fiber) phase $\gamma$ in a way that matches the face-marker geometry.


\subsection{Bloch vector and the SU(2)\texorpdfstring{\(\to\)}{->}SO(3) map}
\label{sec:bloch_vector}

Let \(\ssigma_x,\ssigma_y,\ssigma_z\) denote the Pauli matrices.
The (unit) Bloch vector associated with \(\ssp\) is defined by
\begin{equation}
\label{eq:bloch_def}
\nn(\ssp)
:=
\frac{1}{\ssp^\dagger\ssp}
\begin{bmatrix}
\ssp^\dagger\ssigma_x\ssp\\
\ssp^\dagger\ssigma_y\ssp\\
\ssp^\dagger\ssigma_z\ssp
\end{bmatrix}
\in\Rt,
\qquad
\norm{\nn(\ssp)}=1,
\end{equation}
and for the half-angle parametrization \eqref{eq:spinor_symmetric_halfangle} one obtains the standard expression
\begin{equation}
\label{eq:bloch_angles}
\nn(\ssp)=
\begin{bmatrix}
\sin\theta\cos\phi\\
\sin\theta\sin\phi\\
\cos\theta
\end{bmatrix}.
\end{equation}
This is the usual SU(2) double-cover map to SO(3) (via the adjoint action on Pauli matrices); see \cite{altmann2013rotations,penrose1984spinors}.

In the present visualization, the \emph{primary datum} is the rotated equatorial plane \(E'\) (Section~\ref{sec:hyperchords_setup}), with chosen unit normal \(\vv\).
We identify this normal with the Bloch direction,
\begin{equation}
\label{eq:v_equals_bloch}
\vv=\nn(\ssp),
\end{equation}
so that the polar angle \(\theta\) appearing in Section~\ref{sec:hyperchords_setup} (defined by \(\vv\cdot\eez=\cos\theta\)) is exactly the Bloch polar angle in \eqref{eq:bloch_angles}. (Choosing the opposite normal $-\vv$ corresponds to $\theta\mapsto\pi-\theta$ together with the exchange of the two components; the construction is otherwise unchanged.)

The azimuth \(\phi\) is read geometrically from the equator via the tangency points \(P_\uparrow\) and \(P_\downarrow\) (Section~\ref{sec:hyperchords_equator}).

\subsection{Born-rule weights from face areas}
\label{sec:reading_probabilities}

From Section~\ref{sec:hyperchords_faces}, the two plane-dihedral circles have radii
\[
r_\uparrow=u\cos(\theta/2),
\qquad
r_\downarrow=u\sin(\theta/2).
\]
Hence the corresponding disk areas
\[
A_\uparrow=\pi r_\uparrow^2=\pi u^2\cos^2(\theta/2),
\qquad
A_\downarrow=\pi r_\downarrow^2=\pi u^2\sin^2(\theta/2)
\]
satisfy the normalized relations
\begin{equation}
\label{eq:born_areas}
\frac{A_\uparrow}{\pi u^2}=\cos^2(\theta/2)=\frac{\|\spu\|^2}{u},
\qquad
\frac{A_\downarrow}{\pi u^2}=\sin^2(\theta/2)=\frac{\|\spd\|^2}{u}.
\end{equation}
Thus, at fixed \(u\), the plane-dihedral face areas encode the standard two-level Born weights.

\subsection{Equivalent encoding via minimal projected lengths}
\label{sec:born_projected_lengths}

A second (independent) encoding comes from \emph{minimal} orthogonal projections onto the \xyplane.
The clean way to see it is by slicing the sphere with vertical planes parallel to the meridian plane determined by the common point \(Q\) and the tangency points \(P_\uparrow,P_\downarrow\) (cf. Figure~\ref{fig:hyperchords_schematic}).

Fix such a vertical slice at signed distance \(|y_d|\le u\) from the origin.
Its intersection with the sphere is a circle of radius
\[
r_d=\sqrt{u^2-y_d^2}.
\]
In that slice, the two relevant equator points are \((\pm r_d,0)\) (in \((x,z)\)-coordinates), while the highest point of the rotated equator occurs at
\[
H'(y_d)=(-r_d\cos\theta,\; r_d\sin\theta),
\]
(where, without loss of generality, we have rotated azimuth so that the meridian direction coincides with the \(x\)-axis; the final formulas are azimuth-independent).

Now consider the two Euclidean segments in the slice connecting \(H'(y_d)\) to the two equator points \((\pm r_d,0)\).
Their orthogonal projections onto the \xyplane coincide with their \(x\)-projections in the slice, and therefore have lengths
\begin{equation}
\label{eq:slice_lengths}
\ell_+(y_d)=r_d(1+\cos\theta),
\qquad
\ell_-(y_d)=r_d(1-\cos\theta).
\end{equation}
At \(y_d=0\) this yields the extremal projected chord lengths
\[
\ell_+(0)=u(1+\cos\theta)=2u\cos^2(\theta/2),
\qquad
\ell_-(0)=u(1-\cos\theta)=2u\sin^2(\theta/2),
\]
so that the normalized ratio reproduces the Born weights:
\begin{equation}
\label{eq:born_from_lengths}
\frac{\ell_+(0)}{\ell_+(0)+\ell_-(0)}=\cos^2(\theta/2),
\qquad
\frac{\ell_-(0)}{\ell_+(0)+\ell_-(0)}=\sin^2(\theta/2).
\end{equation}

\subsection{Equivalent encoding via projected areas (slicing proof)}
\label{sec:born_projected_areas}

Integrating the slice widths \eqref{eq:slice_lengths} over all vertical slices gives the projected areas on the \xyplane{} of the two corresponding surface patches (equivalently: the projected areas obtained by stacking the slice projections):
\begin{align}
\Pi_+
&=\int_{-u}^{u}\ell_+(y_d)\,dy_d
=(1+\cos\theta)\int_{-u}^{u}\sqrt{u^2-y_d^2}\,dy_d
=(1+\cos\theta)\frac{\pi u^2}{2}
=\pi u^2\cos^2(\theta/2),
\\
\Pi_-
&=\int_{-u}^{u}\ell_-(y_d)\,dy_d
=(1-\cos\theta)\frac{\pi u^2}{2}
=\pi u^2\sin^2(\theta/2).
\end{align}
Hence
\begin{equation}
\label{eq:born_from_projected_areas}
\frac{\Pi_+}{\Pi_++\Pi_-}=\cos^2(\theta/2),
\qquad
\frac{\Pi_-}{\Pi_++\Pi_-}=\sin^2(\theta/2),
\end{equation}
showing that \emph{disk areas}, \emph{minimal projected lengths}, and \emph{projected slice-integrated areas} all encode the same Born weights.

\subsection{Geometric meaning of \texorpdfstring{\(\phi\)}{phi} and \texorpdfstring{\(\gamma\)}{gamma}}
\label{sec:phase_geometry}

Finally, the remaining angular degrees of freedom in \eqref{eq:spinor_symmetric_halfangle} are visible as in-face rotations of the two plane-dihedral faces.
Using \eqref{eq:phase_split_defs_secIV}, the relative phase $\phi$ is the difference of component arguments,
$\phi=\psi_\downarrow-\psi_\uparrow$, while the fiber (global) phase $\gamma$ is their average,
$\gamma=\tfrac12(\psi_\uparrow+\psi_\downarrow)$.

As summarized in Section~\ref{sec:hyperchords_phase}, the distinguished diameter directions drawn on $c_\uparrow$ and $c_\downarrow$ rotate within their respective face planes by angles determined by $\psi_\uparrow$ and $\psi_\downarrow$.
Therefore:
\begin{itemize}
\item the \emph{relative} twist between the two face markers encodes $\phi$,
\item a \emph{common} twist of both markers encodes $\gamma$ (i.e.\ the choice of representative along the $U(1)$ Hopf fiber).
\end{itemize}
In contrast, bilinear quantities such as the Bloch vector $\nn(\ssp)$ are insensitive to $\gamma$ and depend only on $(\theta,\phi)$.
